\begin{statementbox}
	\textbf{\textcolor{CStatement}{解决什么问题？}}

	看到对数中的乘、除、乘方，或者遇到不同底数的对数，怎样少算几步？把它们与学过的指数运算联系起来，就有下面四条规律。

	\medskip
	\textbf{\textcolor{CStatement}{使用条件}}
	\begin{description}[style=nextline, leftmargin=2.8em, labelsep=0.5em, itemsep=0.35em]
		\item[\textbf{先检查：}] 底数 $a,c>0$ 且 $a,c\neq1$；真数 $M,N,b>0$；指数 $m$ 可以是任意实数。
	\end{description}

	\textbf{\textcolor{CStatement}{四条结论}}
	\begin{description}[style=nextline, leftmargin=2.8em, labelsep=0.5em, itemsep=0.45em]
		\item[\textbf{结论一（积变加）：}]
			同底对数里，真数相乘，可以改为对数相加。
			\[
				\log_a(MN)=\log_a M+\log_a N.
			\]
			例如 $\log_2(8\times4)=3+2=5$。

		\item[\textbf{结论二（商变减）：}]
			同底对数里，真数相除，可以改为对数相减。
			\[
				\log_a\frac{M}{N}=\log_a M-\log_a N.
			\]
			例如 $\log_2(8\div4)=3-2=1$。

		\item[\textbf{结论三（指数提前）：}]
			真数有乘方，可以把指数移到对数前面相乘。
			\[
				\log_a M^m=m\log_a M.
			\]
			例如 $\log_2(8^2)=2\times3=6$。

		\item[\textbf{结论四（换底公式）：}]
			底数不好算时，换成一个熟悉的底数 $c$：\textbf{上面放原真数，下面放原底数}。
			\[
				\log_a b=\frac{\log_c b}{\log_c a}.
			\]
			例如 $\log_4 8=\dfrac{\log_2 8}{\log_2 4}=\dfrac32$。
	\end{description}
\end{statementbox}
